\documentclass[11pt,a4paper]{article}
\usepackage[T1]{fontenc}
\usepackage[utf8]{inputenc}
\usepackage{lmodern}
\usepackage{amsmath,amssymb,amsthm,mathtools}
\usepackage{microtype}
\usepackage[margin=27mm]{geometry}
\usepackage{enumitem}
\usepackage{xcolor}
\usepackage[colorlinks=true,linkcolor=blue!50!black,citecolor=blue!50!black,urlcolor=blue!50!black]{hyperref}
\hypersetup{pdftitle={Local Edge Obstructions and Beta-Prime Classification},pdfauthor={Keio}}
\setlength{\emergencystretch}{2em}
\setlist[enumerate]{itemsep=3pt,topsep=5pt}
\newtheorem{theorem}{Theorem}[section]
\newtheorem{proposition}[theorem]{Proposition}
\newtheorem{lemma}[theorem]{Lemma}
\newtheorem{corollary}[theorem]{Corollary}
\theoremstyle{definition}
\newtheorem{definition}[theorem]{Definition}
\newtheorem{example}[theorem]{Example}
\theoremstyle{remark}
\newtheorem{remark}[theorem]{Remark}
\newcommand{\R}{\mathbb R}
\newcommand{\C}{\mathbb C}
\newcommand{\Cp}{\mathbb C^+}
\newcommand{\Cm}{\mathbb C^-}
\newcommand{\FID}{\mathrm{FID}}
\newcommand{\Bad}{\left(\tfrac12,\tfrac32\right)}
\newcommand{\Good}{\left(0,\tfrac12\right]\cup\left[\tfrac32,\infty\right)}
\DeclareMathOperator{\supp}{supp}
\DeclareMathOperator{\Imn}{Im}
\DeclareMathOperator{\Ren}{Re}
\title{Local Edge Obstructions to Free Infinite Divisibility\\
and the Classical Beta-Prime Classification}
\author{Keio}
\date{15 November 2026\\\small Research manuscript}

\begin{document}
\maketitle

\begin{abstract}
We prove that a probability measure cannot be freely infinitely
divisible if, at a point with a one-sided gap and no atom, its density
is asymptotic to $ct^{p-1}$ on the occupied side, where $c>0$ and
$1/2<p<3/2$. The measure away from the edge is arbitrary, and no
moment or complex-analytic regularity assumptions are imposed.
The proof converts the Nevanlinna representation of the Voiculescu
transform into a third-derivative sign condition for a real inverse
of the reciprocal Cauchy transform. An exact expression in four
real Cauchy integrals violates this condition throughout the stated
interval. Four asymptotic cases include the logarithmic exponent
$p=1$ and the possible cancellation of the finite boundary Cauchy
value. Together with Hasebe's established sufficient results, the
local theorem gives the complete classical beta-prime classification:
$\beta'_{p,q}$ is freely infinitely divisible exactly when
$p\in(0,1/2]\cup[3/2,\infty)$, for every $q>0$.
It also gives the corresponding gamma classification, an exclusion
region for beta distributions, and the exact criterion for scaled
Fisher $F$ laws. The endpoints are sharp.
\end{abstract}

\noindent\textbf{Status.} The arguments are presented as research proofs.
External peer review and discovery priority have not been established.

\section{The local theorem and its consequences}

For Borel probability measures on $\R$, let $\boxplus$ denote free
additive convolution. A probability measure $\mu$ is freely infinitely
divisible, abbreviated FID, if for every integer $n\ge1$ there is a
probability measure $\mu_n$ such that $\mu=\mu_n^{\boxplus n}$.
All divisibility assertions below concern free additive convolution.
The beta, beta-prime, and gamma laws under discussion are the
classical probability distributions with the explicitly stated
densities.

\begin{theorem}[Local one-sided edge obstruction]\label{thm:edge}
Let $\mu$ be a Borel probability measure on $\R$. Suppose that there
exist $x_0\in\R$, $\delta>0$, $c>0$, and $p\in\Bad$ such that
\begin{enumerate}[label=\textup{(\roman*)}]
\item $\mu((x_0-\delta,x_0))=0$ and $\mu(\{x_0\})=0$;
\item the restriction of $\mu$ to $(x_0,x_0+\delta)$ has a density
$r$, with a version satisfying
\begin{equation}\label{eq:edge-assumption}
 r(x_0+t)=ct^{p-1}(1+o(1)),\qquad t\downarrow0.
\end{equation}
\end{enumerate}
Then $\mu$ is not FID. No restriction is imposed on $\mu$ outside
$(x_0-\delta,x_0+\delta)$, and no finite moments are required.
\end{theorem}

Reflection gives the same conclusion when the occupied and empty
sides are reversed. Translation, reflection, and positive scaling
preserve free infinite divisibility: this follows from their
compatibility with sums of freely independent variables, and the
inverse affine map gives the converse. These facts will also be used
in the distributional applications.

Theorem~\ref{thm:edge} answers the local extension question posed as
Problem~5.7 in Hasebe's original 2013 manuscript
\cite{Hasebe2013} and retained in the published paper
\cite{Hasebe2014}. It covers the entire interval $(1/2,3/2)$.
It also replaces the analytic continuation and sector-uniform
remainder hypotheses of Hasebe's Theorem~5.1 by the real asymptotic
\eqref{eq:edge-assumption}. The proof uses real integral asymptotics;
it does not analytically continue the density.

For orientation, the principal consequence is as follows. The full
proof and the other applications appear in Section~\ref{sec:applications}.

\begin{theorem}[Classical beta-prime classification]\label{thm:bp-main}
For every $p,q>0$, let $\beta'_{p,q}$ have density
\begin{equation}\label{eq:bp-density}
 \frac{x^{p-1}}{B(p,q)(1+x)^{p+q}}\,\mathbf1_{(0,\infty)}(x),
 \qquad B(p,q)=\frac{\Gamma(p)\Gamma(q)}{\Gamma(p+q)}.
\end{equation}
Then
\begin{equation}\label{eq:bp-classification}
 \beta'_{p,q}\text{ is FID}
 \quad\Longleftrightarrow\quad
 p\in\Good.
\end{equation}
\end{theorem}

The sufficient direction of \eqref{eq:bp-classification} is Hasebe's
published Theorem~1.2(3), proved more precisely in Theorem~4.7
\cite{Hasebe2014}. The local theorem supplies the complementary
necessary direction. Thus the classification is independent of the
second shape parameter $q$.

\section{A necessary inverse-derivative sign}

\subsection{The analytic characterization and a global identity}

For a probability measure $\mu$, define its Cauchy transform and
reciprocal Cauchy transform by
\begin{equation}\label{eq:G-F}
 G_\mu(z)=\int_\R\frac{1}{z-t}\,\mu(dt),\qquad
 F_\mu(z)=\frac1{G_\mu(z)},\qquad z\in\Cp.
\end{equation}
The transform $G_\mu$ takes values in $\Cm$, so it has no zero in
$\Cp$, and $F_\mu$ maps $\Cp$ into $\Cp$. A compositional inverse
of $F_\mu$ exists in a suitable truncated cone at infinity. The
Voiculescu transform is the inverse germ minus the identity:
\[
 \phi_\mu(w)=F_\mu^{-1}(w)-w.
\]
The Bercovici--Voiculescu analytic characterization says that $\mu$
is FID if and only if this germ extends analytically to $\Cp$ with
\begin{equation}\label{eq:pick}
 \Imn\phi_\mu(w)\le0\qquad(w\in\Cp).
\end{equation}
We use the formulation in Hasebe~\cite[Theorem~2.4]{Hasebe2014};
the original general theorem is in \cite{BV1993}.

Assume temporarily that $\mu$ is FID, and put
\begin{equation}\label{eq:H}
 H(w)=w+\phi_\mu(w),\qquad w\in\Cp.
\end{equation}
The extension $H$ need not be assumed to be univalent on all of
$\Cp$. What is needed, and what follows from analyticity, is the
identity
\begin{equation}\label{eq:global-inverse}
 H(F_\mu(z))=z\qquad(z\in\Cp).
\end{equation}
Indeed, it holds on a nonempty open region near infinity by the
definition of the inverse germ. Both sides are analytic throughout
$\Cp$, because $F_\mu(\Cp)\subset\Cp$. The identity theorem gives
\eqref{eq:global-inverse} on that connected domain. This step
identifies $H$ with the physical inverse branch at every real gap
point used below.

\subsection{Real gaps and the Nevanlinna representation}

The anti-Pick function in \eqref{eq:pick} has a Nevanlinna
representation
\begin{equation}\label{eq:nevanlinna}
 \phi_\mu(w)
 =a-bw+\int_\R\frac{1+tw}{w-t}\,\sigma(dt),
 \qquad a\in\R,\ b\ge0,
\end{equation}
where $\sigma$ is a finite positive measure. For the normalized
Voiculescu transform the linear coefficient is zero; allowing the
term $-bw$ is harmless because its third derivative vanishes.
No moment condition on $\mu$ is needed for this representation.

\begin{lemma}[Third derivative at a real inverse point]\label{lem:sign}
Suppose $\mu$ is FID. Let $x\in\R\setminus\supp\mu$ and suppose
$G_\mu(x)\ne0$. Write $w=F_\mu(x)$. Then $H$ extends analytically
and real-valuedly through a neighborhood of $w$, agrees there with
the local inverse of $F_\mu$, and satisfies
\begin{equation}\label{eq:H-third-sign}
 H'''(w)
 =-6\int_\R\frac{1+t^2}{(w-t)^4}\,\sigma(dt)\le0.
\end{equation}
\end{lemma}

\begin{proof}
The Cauchy transform is analytic near $x$ and real on a real
neighborhood of $x$. Differentiation gives
\[
 G_\mu'(x)=-\int_\R\frac{1}{(x-t)^2}\,\mu(dt)<0,
 \qquad
 F_\mu'(x)
 =\frac{\displaystyle\int_\R(x-t)^{-2}\,\mu(dt)}{G_\mu(x)^2}>0.
\]
The analytic inverse function theorem supplies a local inverse $J$
at $w$. It is real on a real interval and has positive derivative
there after the neighborhoods are made smaller. Consequently it
maps nearby points in the upper half-plane into the upper
half-plane. For such points $\zeta$, \eqref{eq:global-inverse} gives
\[
 H(\zeta)=H(F_\mu(J(\zeta)))=J(\zeta).
\]
Thus $J$ provides the asserted real analytic continuation of $H$,
and hence also of $\phi_\mu$.

Let $J$ be a compact real interval inside this continuation
neighborhood, with $w$ in its interior. The representation gives
\[
 -\Imn\phi_\mu(u+iv)
 =bv+v\int_\R\frac{1+t^2}{(u-t)^2+v^2}\,\sigma(dt),
 \qquad u\in J,\quad v>0.
\]
Real analytic continuation implies that the left side is $O(v)$
uniformly on $J$. Integrating in $u$ shows that the nonnegative
integral term has integral $O(v)$ over $J$. For each $t$ in the
interior $J^\circ$, its inner Poisson-kernel integral tends to
$\pi(1+t^2)$. Tonelli's theorem and Fatou's lemma therefore give
\[
 \pi\int_{J^\circ}(1+t^2)\,\sigma(dt)\le0.
\]
Thus $\sigma(J^\circ)=0$. This is the real-gap consequence of
Stieltjes inversion, with no assumption about absolute continuity
of $\sigma$.

It is now legitimate to differentiate \eqref{eq:nevanlinna} three
times at $w$. On the support of $\sigma$ the denominator stays away
from zero locally, while the differentiated integrand is
$O((1+|t|)^{-2})$ at infinity. Positivity and finiteness of $\sigma$
justify differentiation under the integral. Since
\[
 \frac{d^3}{dw^3}\frac{1+tw}{w-t}
 =-\frac{6(1+t^2)}{(w-t)^4},
\]
we obtain \eqref{eq:H-third-sign}.
\end{proof}

\begin{remark}
The local inverse conclusion is essential. An arbitrary analytic
continuation of an inverse germ cannot simply be identified with a
chosen branch on a real gap. Here that identification follows from
\eqref{eq:global-inverse} and the positivity of $F_\mu'(x)$.
\end{remark}

\section{An exact certificate in four real integrals}

Translate the endpoint in Theorem~\ref{thm:edge} to $x_0=0$; affine
invariance permits this without changing the question of FID.
For $0<s<\delta/2$, define
\begin{equation}\label{eq:Ak}
 A_k(s)=\int_\R\frac{1}{(t+s)^k}\,\mu(dt),
 \qquad k=1,2,3,4.
\end{equation}
These integrals are finite because $-s$ lies in the open gap
$(-\delta,0)$. The odd-indexed integrals can be signed: the measure
may have mass to the left of $-\delta$. In contrast, $A_2(s)>0$.
They are smooth in $s$, with
\begin{equation}\label{eq:Ak-derivative}
 A_k'(s)=-kA_{k+1}(s),\qquad k=1,2,3.
\end{equation}

Whenever $A_1(s)\ne0$, put
\begin{equation}\label{eq:f}
 f(s)=\frac1{A_1(s)}.
\end{equation}
Then $G_\mu(-s)=-A_1(s)$ and
\begin{equation}\label{eq:F-f}
 F_\mu(-s)=-f(s).
\end{equation}
The derivative $f'(s)=A_2(s)/A_1(s)^2$ is strictly positive.

\begin{proposition}[Exact inverse certificate]\label{prop:N}
At a gap point for which $A_1(s)\ne0$, the third derivative of the
local inverse of $F_\mu$ is
\begin{equation}\label{eq:exact-H-third}
 H'''(-f(s))
 =\frac{6A_1(s)^4}{A_2(s)^5}\,N(s),
\end{equation}
where
\begin{equation}\label{eq:N}
 N(s)=A_2^4-2A_1A_2^2A_3
       +A_1^2\bigl(2A_3^2-A_2A_4\bigr).
\end{equation}
Here the notation $H$ refers to that local inverse; under the FID
assumption it agrees with \eqref{eq:H}. Consequently, FID requires
\begin{equation}\label{eq:N-necessary}
 N(s)\le0
\end{equation}
at every such gap point.
\end{proposition}

\begin{proof}
Using \eqref{eq:Ak-derivative}, direct differentiation gives
\begin{align}
 f'&=\frac{A_2}{A_1^2},\notag\\
 f''&=\frac{2(A_2^2-A_1A_3)}{A_1^3},\label{eq:f-derivatives}\\
 f'''&=\frac{6(A_1^2A_4-2A_1A_2A_3+A_2^3)}{A_1^4}.\notag
\end{align}
Let $g$ be the local inverse of $f$. Equation~\eqref{eq:F-f} gives
$H(w)=-g(-w)$, and hence $H'''(w)=g'''(-w)$. Inverse-function
differentiation yields
\[
 g'''(f(s))=\frac{3(f''(s))^2-f'''(s)f'(s)}{(f'(s))^5}.
\]
Inserting \eqref{eq:f-derivatives} and collecting terms gives
\eqref{eq:exact-H-third}--\eqref{eq:N}. Its prefactor is strictly
positive. Lemma~\ref{lem:sign} therefore proves
\eqref{eq:N-necessary} under the FID assumption.
\end{proof}

The proof of Theorem~\ref{thm:edge} will show that $A_1(s)$ is
nonzero and $N(s)$ is positive for all sufficiently small $s>0$.
This gives a local contradiction without requiring finite moments,
free cumulants, or Hankel determinants.

\section{Edge asymptotics under a real density hypothesis}

We record the estimates needed to evaluate \eqref{eq:N}. In this
section $r(t)=ct^{p-1}(1+o(1))$ as $t\downarrow0$ and the local
gap and absence of an atom are as in Theorem~\ref{thm:edge}.
The beta integral is
\[
 B(u,v)=\int_0^\infty\frac{x^{u-1}}{(1+x)^{u+v}}\,dx
       =\frac{\Gamma(u)\Gamma(v)}{\Gamma(u+v)},\qquad u,v>0.
\]

\begin{lemma}[Power and logarithmic estimates]\label{lem:asymptotics}
If $k>p$, then
\begin{equation}\label{eq:Ak-asymptotic}
 A_k(s)\sim cB(p,k-p)s^{p-k}
 =c\frac{\Gamma(p)\Gamma(k-p)}{\Gamma(k)}s^{p-k}.
\end{equation}
For $p=1$ one also has
\begin{equation}\label{eq:A1-log}
 A_1(s)\sim c\log(1/s).
\end{equation}
For $1<p<2$, the finite limit
\begin{equation}\label{eq:a}
 a=A_1(0)=\int_\R\frac1t\,\mu(dt)
\end{equation}
exists as an absolutely convergent integral, and
\begin{equation}\label{eq:A1-difference}
 A_1(s)-a
 \sim-cB(p-1,2-p)s^{p-1}.
\end{equation}
\end{lemma}

\begin{proof}
The measure on $\R\setminus(-\delta,\delta)$ contributes $O(1)$
to each $A_k(s)$, uniformly for $0<s<\delta/2$. This uses only its
finite total mass: $|t+s|\ge\delta/2$ there. The gap contributes
nothing, and there is no atom at zero.

Fix $\varepsilon\in(0,\delta)$ sufficiently small. On
$(0,\varepsilon)$, the density is bounded above and below by
arbitrarily close positive multiples of $ct^{p-1}$. After the
substitution $t=su$, its comparison integral is
\[
 \int_0^\varepsilon\frac{ct^{p-1}}{(t+s)^k}\,dt
 =cs^{p-k}\int_0^{\varepsilon/s}
            \frac{u^{p-1}}{(1+u)^k}\,du.
\]
If $k>p$, the last integrand is integrable at zero and infinity,
and its integral tends to $B(p,k-p)$. The contribution from
$(\varepsilon,\delta)$ is bounded, even if the density there has
no further regularity. Since $s^{p-k}\to\infty$, these bounded
terms are negligible. Letting the comparison error tend to zero
proves \eqref{eq:Ak-asymptotic}.

When $p=1$, the corresponding comparison integral for $A_1$ is
$c\log((s+\varepsilon)/s)$. The same bounds give
\eqref{eq:A1-log}.

If $1<p<2$, integrability of $r(t)/t$ at zero proves that
\eqref{eq:a} is absolutely convergent; outside the local interval,
$|t|$ is bounded away from zero. The exact difference formula is
\begin{equation}\label{eq:difference-exact}
 A_1(s)-a=-s\int_\R\frac{1}{t(t+s)}\,\mu(dt).
\end{equation}
The exterior contribution is $O(s)$. On $(0,\varepsilon)$,
comparison and scaling give the leading term
\[
 -cs^{p-1}\int_0^{\varepsilon/s}
              \frac{u^{p-2}}{1+u}\,du
 \sim-cB(p-1,2-p)s^{p-1}.
\]
The remaining local interval contributes $O(s)$ as well. Since
$p-1<1$, those terms are lower order. This proves
\eqref{eq:A1-difference}.
\end{proof}

\begin{remark}
The comparison proof requires only the real asymptotic
\eqref{eq:edge-assumption}. It does not differentiate the density
or assume a convergent expansion for it. Mass on the remote
negative half-line affects the value of $a$ in \eqref{eq:a}, but
not the divergent coefficients in \eqref{eq:Ak-asymptotic}.
\end{remark}

\section{The four cases proving the obstruction}\label{sec:fourcases}

Assume for contradiction that the measure in Theorem~\ref{thm:edge}
is FID. By Proposition~\ref{prop:N}, it is enough to prove eventual
positivity of $N(s)$ and nonvanishing of $A_1(s)$. All asymptotic
relations in this section are taken as $s\downarrow0$.

\subsection{Case I: \texorpdfstring{$1/2<p<1$}{1/2 < p < 1}}

Put
\[
 r_0=1-p\in(0,1/2),\qquad
 D=c\Gamma(p)\Gamma(1-p)>0.
\]
To avoid confusion with the density $r$, the parameter in the next
formulas is denoted $r_0$. By \eqref{eq:Ak-asymptotic},
\begin{equation}\label{eq:pochhammer}
 A_k(s)\sim D\frac{(r_0)_{k-1}}{(k-1)!}s^{p-k},
 \qquad k=1,2,3,4,
\end{equation}
where $(r_0)_j=r_0(r_0+1)\cdots(r_0+j-1)$ and $(r_0)_0=1$.
Every monomial in \eqref{eq:N} has the same power $s^{4p-8}$.
Its coefficient, after division by $D^4$, is
\begin{align*}
 &r_0^4-r_0^3(r_0+1)
   +\frac{r_0^2(r_0+1)^2}{2}
   -\frac{r_0^2(r_0+1)(r_0+2)}6\\
 &\hspace{28mm}=
 \frac{r_0^2(1-r_0)(1-2r_0)}6.
\end{align*}
Thus
\begin{equation}\label{eq:N-case-I}
 N(s)\sim
 \frac{D^4r_0^2(1-r_0)(1-2r_0)}6\,s^{4p-8}>0.
\end{equation}
Here and below a positive asymptotic means eventual strict
positivity. Since $A_1(s)\sim Ds^{p-1}>0$, the necessary inequality
\eqref{eq:N-necessary} is contradicted.

\subsection{Case II: \texorpdfstring{$p=1$}{p = 1}}

The logarithmic estimate and the power estimates give
\begin{equation}\label{eq:case-II-A}
 A_1(s)\sim c\log(1/s),\quad
 A_2(s)\sim\frac cs,\quad
 A_3(s)\sim\frac{c}{2s^2},\quad
 A_4(s)\sim\frac{c}{3s^3}.
\end{equation}
Consequently,
\[
 2A_3(s)^2-A_2(s)A_4(s)\sim\frac{c^2}{6s^4}.
\]
In \eqref{eq:N}, the term containing $A_1^2$ has a
logarithm-squared factor; the other two terms have at most one
logarithm. Hence
\begin{equation}\label{eq:N-case-II}
 N(s)\sim\frac{c^4}{6s^4}\bigl(\log(1/s)\bigr)^2>0.
\end{equation}
The integral $A_1(s)$ is positive for all sufficiently small $s$,
so this again contradicts \eqref{eq:N-necessary}. The exponent
$p=1$ is therefore included, without a limiting argument from
the adjacent cases.

\subsection{Case III: \texorpdfstring{$1<p<3/2$ and $A_1(0)\ne0$}{1 < p < 3/2 and A1(0) nonzero}}

Let $a=A_1(0)\ne0$, and set
\[
 d_2=c\Gamma(p)\Gamma(2-p)>0.
\]
Lemma~\ref{lem:asymptotics} gives
\begin{align}
 A_1(s)&\longrightarrow a,\notag\\
 A_2(s)&\sim d_2s^{p-2},\notag\\
 A_3(s)&\sim\frac{d_2(2-p)}2s^{p-3},\label{eq:case-III-A}\\
 A_4(s)&\sim\frac{d_2(2-p)(3-p)}6s^{p-4}.\notag
\end{align}
It follows that
\[
 2A_3^2-A_2A_4
 \sim\frac{d_2^2(2-p)(3-2p)}6s^{2p-6}.
\]
This leading coefficient is positive. The term $A_2^4$ has
relative order $s^{2p-2}$ and the term $A_1A_2^2A_3$ has relative
order $s^{p-1}$ compared with $A_1^2(2A_3^2-A_2A_4)$. Both tend
to zero because $p>1$. Thus
\begin{equation}\label{eq:N-case-III}
 N(s)\sim
 \frac{a^2d_2^2(2-p)(3-2p)}6s^{2p-6}>0.
\end{equation}
The sign of $a$ is irrelevant, and $A_1(s)$ has that nonzero
sign eventually. This yields the required contradiction.

\subsection{Case IV: \texorpdfstring{$1<p<3/2$ and $A_1(0)=0$}{1 < p < 3/2 and A1(0) = 0}}

This possibility must be retained for arbitrary measures on $\R$:
remote mass to the left of the gap can cancel the finite boundary
integral. Formula~\eqref{eq:A1-difference} now gives
\begin{equation}\label{eq:A1-zero-case}
 A_1(s)\sim
 -c\Gamma(p-1)\Gamma(2-p)s^{p-1}
 =Ds^{p-1},
 \qquad D=c\Gamma(p)\Gamma(1-p)<0.
\end{equation}
Put $r_0=1-p\in(-1/2,0)$. Equations
\eqref{eq:Ak-asymptotic} and \eqref{eq:A1-zero-case} show that
the same four formulas \eqref{eq:pochhammer} hold. Although $D$
and $r_0$ are negative, their combinations for $k\ge2$ are positive,
in agreement with the local integral asymptotics.

The coefficient computation in Case I is algebraic and remains
valid for this value of $r_0$. Therefore
\begin{equation}\label{eq:N-case-IV}
 N(s)\sim
 \frac{D^4r_0^2(1-r_0)(1-2r_0)}6s^{4p-8}>0,
\end{equation}
since all displayed factors apart from even powers are positive.
Also $A_1(s)<0$ for sufficiently small $s$, so it is nonzero.
This contradicts \eqref{eq:N-necessary} in the final case.

\begin{proof}[Completion of Theorem~\ref{thm:edge}]
The four cases exhaust $1/2<p<3/2$, including both possibilities
for the finite boundary integral when $p>1$. In each case,
$A_1(s)\ne0$ and $N(s)>0$ for all sufficiently small positive $s$.
Proposition~\ref{prop:N} and Lemma~\ref{lem:sign} exclude FID.
Translation returns the conclusion to the original endpoint $x_0$.
\end{proof}

\section{Classifications and distributional applications}
\label{sec:applications}

\subsection{Classical beta-prime laws}

\begin{proof}[Proof of Theorem~\ref{thm:bp-main}]
For $p,q>0$, the density \eqref{eq:bp-density} is zero to the left
of zero, has no atom at zero, and satisfies
\[
 \frac{x^{p-1}}{B(p,q)(1+x)^{p+q}}
 \sim\frac1{B(p,q)}x^{p-1}\qquad(x\downarrow0).
\]
If $1/2<p<3/2$, Theorem~\ref{thm:edge} excludes FID, for every
$q>0$. If $p\in(0,1/2]\cup[3/2,\infty)$, Hasebe's
Theorem~4.7 establishes membership in the univalent-inverse class
$\mathrm{UI}$, hence FID, for every $q>0$
\cite[Theorem~1.2(3)]{Hasebe2014}. The two ranges exhaust the
positive values of $p$ and include both boundary points in the
sufficient direction.
\end{proof}

This settles the beta-prime assertion in the original
Conjecture~5.6(1) of \cite{Hasebe2013}, and combines it with the
already established sufficient range. Arbitrarily heavy right tails
are permitted: no condition on moments or on the size of $q$ enters
the local obstruction.

\subsection{Gamma laws and sharpness}

\begin{corollary}[Gamma classification]\label{cor:gamma}
Let $\gamma_p$ be the classical gamma law of shape $p>0$ and rate
one, with density
\[
 \frac{x^{p-1}e^{-x}}{\Gamma(p)}\mathbf1_{(0,\infty)}(x).
\]
Then
\[
 \gamma_p\text{ is FID}
 \quad\Longleftrightarrow\quad p\in\Good.
\]
The same classification holds at every positive scale.
\end{corollary}

\begin{proof}
The density has the required edge asymptotic with $c=1/\Gamma(p)$,
so Theorem~\ref{thm:edge} gives non-FID for $p\in(1/2,3/2)$.
Hasebe's sufficient result gives FID on the complementary range
\cite[Corollary~1.3(1)]{Hasebe2014}. Positive scaling preserves
FID in both directions.
\end{proof}

\begin{corollary}[Sharpness of the local exponent interval]
\label{cor:sharpness}
For each $p>0$ outside $(1/2,3/2)$ there is an FID probability
measure with a one-sided atom-free edge of density asymptotic to
$ct^{p-1}$ for some $c>0$. In particular, neither endpoint $1/2$
nor $3/2$ can be added to a universal obstruction based only on
the hypotheses and exponent in Theorem~\ref{thm:edge}.
\end{corollary}

\begin{proof}
Use the gamma laws in the established sufficient range of
Corollary~\ref{cor:gamma}. They satisfy every local support and
density hypothesis at zero.
\end{proof}

The gamma conclusion is an independent consequence of the local
argument, not a claim of first discovery of that classification.
A 2026 Ideosphere preprint already claims a full gamma classification
through finite Hankel certificates \cite{IdeosphereGamma}.
Its version~3, Remark~12.13, explicitly leaves Hasebe's local
Problem~5.7 and Conjecture~5.6(1) outside its scope. This is a
preprint claim whose complete underlying proof has not been
independently audited here. The present argument does not use it.

\subsection{Beta laws: a negative range}

\begin{corollary}[Beta exclusion]\label{cor:beta}
For $p,q>0$, let $\beta_{p,q}$ be the classical beta distribution
with density
\[
 \frac{x^{p-1}(1-x)^{q-1}}{B(p,q)}\mathbf1_{(0,1)}(x).
\]
If $p\in(1/2,3/2)$ or $q\in(1/2,3/2)$, then $\beta_{p,q}$
is not FID.
\end{corollary}

\begin{proof}
At zero the density is asymptotic to $x^{p-1}/B(p,q)$, so the
first exclusion follows from Theorem~\ref{thm:edge}. If $X$ has
law $\beta_{p,q}$, then $1-X$ has law $\beta_{q,p}$. Reflection
and translation preserve FID, so the same argument gives the
second exclusion.
\end{proof}

This proves the negative-range assertion about beta laws in
Conjecture~5.6(1) of \cite{Hasebe2013}. It is not a complete beta
classification. In particular, the positive-region assertion in
Conjecture~5.6(2) is not established by this argument.

\subsection{Scaled Fisher \texorpdfstring{$F$}{F} laws}

For positive real degrees $d_1,d_2$, define the classical Fisher law
$\mathsf F_{d_1,d_2}$ by its beta-prime representation
\begin{equation}\label{eq:Fisher-scaling}
 Y\sim\mathsf F_{d_1,d_2}\quad\Longleftrightarrow\quad
 Y=\frac{d_2}{d_1}X,\qquad
 X\sim\beta'_{d_1/2,d_2/2}.
\end{equation}
This agrees with the usual ratio of independent chi-square
variables divided by their degrees of freedom. Equivalently,
putting $p=d_1/2$ and $q=d_2/2$, its density is
\[
 \frac{(p/q)^p}{B(p,q)}
 x^{p-1}\left(1+\frac pqx\right)^{-p-q}
 \mathbf1_{(0,\infty)}(x).
\]

\begin{corollary}[Fisher classification]\label{cor:Fisher}
For all $d_1,d_2>0$,
\[
 \mathsf F_{d_1,d_2}\text{ is FID}
 \quad\Longleftrightarrow\quad
 d_1\in(0,1]\cup[3,\infty).
\]
The criterion is unchanged by further positive scaling and places
no restriction on $d_2$. For positive integer degrees, the only
excluded numerator degree is $d_1=2$.
\end{corollary}

\begin{proof}
The scale factor $d_2/d_1$ in \eqref{eq:Fisher-scaling} is
strictly positive. Dilation by a positive constant commutes with
free additive convolution, and its inverse dilation does as well.
Thus $Y$ is FID exactly when $X$ is FID. Apply
Theorem~\ref{thm:bp-main} with $p=d_1/2$.
\end{proof}

\section{Relation to the original local problem and limitations}

Hasebe's Theorem~5.1 in the 2013 preprint \cite{Hasebe2013} and
the 2014 publication \cite{Hasebe2014} supplied a local obstruction
for an alternating collection of exponent intervals contained in
$(1/2,3/2)$. Its remainder was required to extend analytically to
appropriate complex sectors and to tend uniformly to zero there.
Problem~5.7 asked whether the exponent restriction could be replaced
by the full interval. Theorem~\ref{thm:edge} proves that extension
with the weaker real asymptotic \eqref{eq:edge-assumption}.

The proof isolates two features that allow the conclusion to be
uniform in the exponent. First, the Pick representation controls
the third derivative of every physical real inverse branch; the
identity \eqref{eq:global-inverse} makes this connection exact.
Second, the rational certificate \eqref{eq:N} has an explicitly
positive leading term in all four cases. The logarithmic case
$p=1$ and the cancellation case $A_1(0)=0$ are both treated
directly, rather than inferred from neighboring parameter values.

The conclusion is local in the original measure. Remote mass may
change the sign of $A_1(0)$ or make it vanish, and may destroy every
positive moment. These possibilities do not affect the proof.
On the other hand, the absence of an atom and the one-sided gap
are hypotheses: no assertion is made here for an edge with a
positive atom, or for a density that is positive on both sides
of the reference point. The beta exclusion is only the range
proved in Corollary~\ref{cor:beta}; the complete classifications
in this paper concern beta-prime, gamma, and the scaled Fisher
laws specified above.

The identities and the four asymptotic coefficient calculations
have been checked independently within the research workflow.
Symbolic checks recovered \eqref{eq:exact-H-third}, the coefficient
$r_0^2(1-r_0)(1-2r_0)/6$, and the identity $N=0$ for a point mass.
Ten high-precision numerical checks, including measures with
$A_1(0)=0$, gave positive $N$ at tested exponents in the excluded
interval. The proof is analytic and does not depend on these finite
checks. Neither the checks nor the references to earlier problems
establish external peer review or discovery priority.

\begin{thebibliography}{9}

\bibitem{BV1993}
H.~Bercovici and D.~Voiculescu,
\emph{Free convolution of measures with unbounded support},
Indiana University Mathematics Journal \textbf{42} (1993), no.~3,
733--773.

\bibitem{Hasebe2013}
T.~Hasebe,
\emph{Free infinite divisibility for beta distributions and related ones},
original preprint, version~1, 4 May 2013,
\href{https://arxiv.org/abs/1305.0924v1}{arXiv:1305.0924v1}.
See Theorem~5.1, Conjecture~5.6, and Problem~5.7.

\bibitem{Hasebe2014}
T.~Hasebe,
\emph{Free infinite divisibility for beta distributions and related ones},
Electronic Journal of Probability \textbf{19} (2014), no.~81, 1--33.
\href{https://doi.org/10.1214/EJP.v19-3448}{doi:10.1214/EJP.v19-3448}.
\href{https://eprints.lib.hokudai.ac.jp/repo/huscap/all/57180/EJP_19_81.pdf}{Author-institution copy}.

\bibitem{IdeosphereGamma}
\emph{Which gamma laws are freely infinitely divisible, and strict
enlargement of the Hankel hierarchy of free cumulants},
Ideosphere preprint, version~3 (2026), especially Remark~12.13,
\href{https://ideosphere.ai/papers/hankel-fid-v3.pdf}{Version~3 manuscript};
\href{https://ideosphere.ai/papers/hankel-fid}{manuscript listing}.
The gamma classification is cited as a preprint claim and is not
used as a proof input here.

\end{thebibliography}
\end{document}
